\subsection{Giới Thiệu}

\textbf{Mục tiêu.} Phần này minh họa toàn bộ quy trình áp dụng các phương pháp data fitting đã
cài đặt ở Phần 1 vào một bộ dữ liệu thực tế. Nhóm không chỉ huấn luyện mô hình hồi quy mà còn xây
dựng pipeline tiền xử lý có thể tái lập, tránh rò rỉ dữ liệu, xử lý missing values, ngoại lai và
đa cộng tuyến trước khi đánh giá mô hình trên tập kiểm tra độc lập.

\textbf{Bài toán.} Với các thông tin quan sát về quỹ đạo, transit, khối lượng hành tinh và đặc
trưng của sao chủ, liệu có thể xây dựng mô hình toán học để dự đoán bán kính hành tinh ngoài Hệ
Mặt Trời hay không? Biến mục tiêu được chọn là:
\begin{equation}
    y = \texttt{pl\_rade},
\end{equation}
trong đó \texttt{pl\_rade} là bán kính hành tinh theo đơn vị bán kính Trái Đất.

\textbf{Quy ước.} Mức ý nghĩa dùng cho các kiểm định và chọn biến thống kê là
$\alpha = 0.05$. Các chỉ số RMSE, MAE và $R^2$ của mô hình tuyến tính chính được tính trên thang
$\log(1+\texttt{pl\_rade})$, vì target đã được log-transform trong pipeline. Khi cần diễn giải
trực tiếp theo bán kính hành tinh, báo cáo có thêm bảng dự báo đã biến đổi ngược về thang gốc.

Toàn bộ kết quả chính thức của Phần 2 được lưu trong thư mục \texttt{part2/real/}. Các file quan
trọng gồm:
\begin{itemize}
    \item \texttt{preprocessing\_metadata.json}: thông tin log-transform, Winsorization, MICE,
          VIF và các feature cuối cùng.
    \item \texttt{model\_results.json}, \texttt{model\_summary.csv}: hệ số, metric train/test và
          kết quả cross-validation của OLS, Ridge, Lasso.
    \item \texttt{advanced\_results.json}, \texttt{kernel\_ridge\_search.csv}: kết quả Kernel
          Ridge Regression với RBF kernel.
    \item Các hình EDA và diagnostic trong \texttt{part2/real/eda/} và
          \texttt{part2/real/diagnostics/}.
\end{itemize}


